\begin{center}
\textbf{Condition d'utilisation d'un médicament}
\end{center}

Le $^{177}$Lu-PSMA est un médicament radiopharmaceutique indiqué dans le traitement
de certaines tumeurs. Il est composé de lutécium 177 ($^{177}$Lu) qui est un isotope
radioactif du lutécium et qui se désintègre par émission $\beta^-$ en hafnium 177 ($^{177}$Hf).

On admet que la fonction $A$ qui modélise l'évolution de l'activité radioactive (en MBq) du lutécium 177 contenu dans une dose de ce médicament, en fonction du temps $t$, exprimé en jour écoulé depuis l'injection, vérifie l'équation différentielle :

\[y' = -0,1y\]

où $y$ est une fonction définie et dérivable sur l'intervalle [0 ;$ +\infty$[ et $y'$ sa dérivée.

\begin{enumerate}[start=7]
\item  Démontrer, en admettant que $A(0) = \np{7400}$, que, pour tout réel $t$ positif ou nul,
on a:
\end{enumerate}

\[A(t)=\np{7400}\e^{-0,1t}\]

L'activité radioactive du lutécium 177 est considérée comme négligeable et sans
risque pour l'environnement lorsqu'elle a perdu plus de 99,9\,\% de sa valeur initiale.
\begin{enumerate}[resume]
\item Montrer que l'activité du lutécium 177 contenu dans cette dose de médicament a
diminué de plus de 99,9\,\% au bout de 70 jours.
\end{enumerate}
Le fabricant du médicament indique dans une notice : \og Stocker dans un sac en
plastique tous les objets non lavables contenant des traces d'urine et placer le sac
dans une poubelle à l'écart. Le jeter aux ordures ménagères au bout de 70 jours\fg.

\begin{enumerate}[resume]
 \item Commenter, en utilisant le résultat de la question précédente, la préconisation du
fabricant.
\end{enumerate}

\bigskip
